% =====================================================================
%  Anatomía de un LLM — Fase II — Módulo 7
%  Optimización
%  Compilar: latexmk -r ../latexmkrc M07-optimizacion.tex
% =====================================================================
\documentclass[11pt,a4paper,twoside,openany]{book}
\usepackage{../estilo-apuntes}

\newcommand{\ApuntesCurso}{Anatomía de un LLM}
\newcommand{\ApuntesModulo}{Módulo 7}

\hypersetup{pdftitle={Módulo 7. Optimización (resumen)},pdfauthor={Anatomía de un LLM}}

% Notación propia del módulo
\newcommand{\T}{^{\top}}
\newcommand{\uno}{\mathbf{1}}
\newcommand{\diag}{\operatorname{diag}}
\newcommand{\Var}{\operatorname{Var}}
\newcommand{\Cov}{\operatorname{Cov}}
\newcommand{\sign}{\operatorname{signo}}
\newcommand{\argmin}{\operatorname*{arg\,min}}
\newcommand{\gb}{\hat g}

\makeatletter
\newlabel{ej:cuadratica}{{1.7}{2}{Funciones cuadráticas}{ejemplo.7}{}}
\newlabel{ej:cuadratica@cref}{{[ejemplo][7][]1.7}{[1][2][]2}}
\newlabel{ej:silla}{{1.13}{4}{Silla}{ejemplo.13}{}}
\newlabel{ej:silla@cref}{{[ejemplo][13][]1.13}{[1][4][]4}}
\newlabel{eje:1-cuartica}{{1.1}{5}{}{ejercicio.1.1}{}}
\newlabel{eje:1-cuartica@cref}{{[ejercicio][1][1]1.1}{[1][4][]5}}
\newlabel{eje:1-silla}{{1.2}{5}{}{ejercicio.1.2}{}}
\newlabel{eje:1-silla@cref}{{[ejercicio][2][1]1.2}{[1][5][]5}}
\newlabel{eje:1-cuadratica}{{1.3}{5}{}{ejercicio.1.3}{}}
\newlabel{eje:1-cuadratica@cref}{{[ejercicio][3][1]1.3}{[1][5][]5}}
\newlabel{eje:1-constantes}{{1.4}{5}{}{ejercicio.1.4}{}}
\newlabel{eje:1-constantes@cref}{{[ejercicio][4][1]1.4}{[1][5][]5}}
\newlabel{eje:1-maximo}{{1.5}{5}{}{ejercicio.1.5}{}}
\newlabel{eje:1-maximo@cref}{{[ejercicio][5][1]1.5}{[1][5][]5}}
\newlabel{eje:1-fuerte}{{1.6}{5}{}{ejercicio.1.6}{}}
\newlabel{eje:1-fuerte@cref}{{[ejercicio][6][1]1.6}{[1][5][]5}}
\newlabel{ej:gd-1d}{{2.7}{7}{Una variable}{ejemplo.7}{}}
\newlabel{ej:gd-1d@cref}{{[ejemplo][7][]2.7}{[1][7][]7}}
\newlabel{ej:gd-2d}{{2.10}{8}{Dos variables}{ejemplo.10}{}}
\newlabel{ej:gd-2d@cref}{{[ejemplo][10][]2.10}{[1][8][]8}}
\newlabel{fig:zigzag}{{2.1}{9}{Descenso de gradiente sobre $f(x,y)=x^{2}+5y^{2}$ con el paso óptimo $\eta ^{\ast }=1/6$, desde $(1,1)$. Cada coordenada se multiplica por $2/3$ (en $x$) o por $-2/3$ (en $y$): la trayectoria zigzaguea entre las paredes del valle. Las elipses son curvas de nivel}{figure.2.1}{}}
\newlabel{fig:zigzag@cref}{{[figure][1][2]2.1}{[1][8][]9}}
\newlabel{ej:rosenbrock}{{2.12}{9}{Función de Rosenbrock}{ejemplo.12}{}}
\newlabel{ej:rosenbrock@cref}{{[ejemplo][12][]2.12}{[1][8][]9}}
\newlabel{eje:2-escala}{{2.1}{10}{}{ejercicio.2.1}{}}
\newlabel{eje:2-escala@cref}{{[ejercicio][1][2]2.1}{[1][10][]10}}
\newlabel{eje:2-1d}{{2.2}{10}{}{ejercicio.2.2}{}}
\newlabel{eje:2-1d@cref}{{[ejercicio][2][2]2.2}{[1][10][]10}}
\newlabel{eje:2-2d}{{2.3}{10}{}{ejercicio.2.3}{}}
\newlabel{eje:2-2d@cref}{{[ejercicio][3][2]2.3}{[1][10][]10}}
\newlabel{eje:2-valle}{{2.4}{10}{}{ejercicio.2.4}{}}
\newlabel{eje:2-valle@cref}{{[ejercicio][4][2]2.4}{[1][10][]10}}
\newlabel{eje:2-paso-largo}{{2.5}{10}{}{ejercicio.2.5}{}}
\newlabel{eje:2-paso-largo@cref}{{[ejercicio][5][2]2.5}{[1][10][]10}}
\newlabel{eje:2-rosenbrock}{{2.6}{10}{}{ejercicio.2.6}{}}
\newlabel{eje:2-rosenbrock@cref}{{[ejercicio][6][2]2.6}{[1][10][]10}}
\newlabel{eje:2-cuadratica}{{2.7}{10}{}{ejercicio.2.7}{}}
\newlabel{eje:2-cuadratica@cref}{{[ejercicio][7][2]2.7}{[1][10][]10}}
\newlabel{ej:minibatch}{{3.4}{12}{}{ejemplo.4}{}}
\newlabel{ej:minibatch@cref}{{[ejemplo][4][]3.4}{[1][12][]12}}
\newlabel{fig:bucle}{{3.1}{12}{Un paso de entrenamiento de un modelo de lenguaje. La retropropagación se estudia en los módulos~9 y~10, el recorte del gradiente en el módulo~11, AdamW en los temas~5 y~6 y el calendario en el tema~6}{figure.3.1}{}}
\newlabel{fig:bucle@cref}{{[figure][1][3]3.1}{[1][12][]12}}
\newlabel{eje:3-motivos}{{3.1}{14}{}{ejercicio.3.1}{}}
\newlabel{eje:3-motivos@cref}{{[ejercicio][1][3]3.1}{[1][14][]14}}
\newlabel{eje:3-minibatch}{{3.2}{14}{}{ejercicio.3.2}{}}
\newlabel{eje:3-minibatch@cref}{{[ejercicio][2][3]3.2}{[1][14][]14}}
\newlabel{eje:3-reemplazo}{{3.3}{14}{}{ejercicio.3.3}{}}
\newlabel{eje:3-reemplazo@cref}{{[ejercicio][3][3]3.3}{[1][14][]14}}
\newlabel{eje:3-suelo}{{3.4}{14}{}{ejercicio.3.4}{}}
\newlabel{eje:3-suelo@cref}{{[ejercicio][4][3]3.4}{[1][14][]14}}
\newlabel{eje:3-sin-reemplazo}{{3.5}{15}{}{ejercicio.3.5}{}}
\newlabel{eje:3-sin-reemplazo@cref}{{[ejercicio][5][3]3.5}{[1][14][]15}}
\newlabel{eje:3-ruido}{{3.6}{15}{}{ejercicio.3.6}{}}
\newlabel{eje:3-ruido@cref}{{[ejercicio][6][3]3.6}{[1][15][]15}}
\newlabel{ej:momentum}{{4.5}{17}{}{ejemplo.5}{}}
\newlabel{ej:momentum@cref}{{[ejemplo][5][]4.5}{[1][17][]17}}
\newlabel{eje:4-efectivo}{{4.1}{19}{}{ejercicio.4.1}{}}
\newlabel{eje:4-efectivo@cref}{{[ejercicio][1][4]4.1}{[1][19][]19}}
\newlabel{eje:4-pasos}{{4.2}{19}{}{ejercicio.4.2}{}}
\newlabel{eje:4-pasos@cref}{{[ejercicio][2][4]4.2}{[1][19][]19}}
\newlabel{eje:4-optimos}{{4.3}{19}{}{ejercicio.4.3}{}}
\newlabel{eje:4-optimos@cref}{{[ejercicio][3][4]4.3}{[1][19][]19}}
\newlabel{eje:4-autovalores}{{4.4}{19}{}{ejercicio.4.4}{}}
\newlabel{eje:4-autovalores@cref}{{[ejercicio][4][4]4.4}{[1][19][]19}}
\newlabel{eje:4-forma}{{4.5}{20}{}{ejercicio.4.5}{}}
\newlabel{eje:4-forma@cref}{{[ejercicio][5][4]4.5}{[1][20][]20}}
\newlabel{eje:4-estabilidad}{{4.6}{20}{}{ejercicio.4.6}{}}
\newlabel{eje:4-estabilidad@cref}{{[ejercicio][6][4]4.6}{[1][20][]20}}
\newlabel{fig:genealogia}{{5.1}{23}{Genealogía de los optimizadores del módulo. AdamW, que se estudia en el \cref {tema:6}, es el optimizador habitual en el preentrenamiento de los modelos de lenguaje abiertos documentados}{figure.5.1}{}}
\newlabel{fig:genealogia@cref}{{[figure][1][5]5.1}{[1][23][]23}}
\newlabel{ej:adam}{{5.9}{23}{}{ejemplo.9}{}}
\newlabel{ej:adam@cref}{{[ejemplo][9][]5.9}{[1][23][]23}}
\newlabel{eje:5-escala-test}{{5.1}{24}{}{ejercicio.5.1}{}}
\newlabel{eje:5-escala-test@cref}{{[ejercicio][1][5]5.1}{[1][24][]24}}
\newlabel{eje:5-rmsprop}{{5.2}{24}{}{ejercicio.5.2}{}}
\newlabel{eje:5-rmsprop@cref}{{[ejercicio][2][5]5.2}{[1][24][]24}}
\newlabel{eje:5-adam}{{5.3}{24}{}{ejercicio.5.3}{}}
\newlabel{eje:5-adam@cref}{{[ejercicio][3][5]5.3}{[1][24][]24}}
\newlabel{eje:5-precond}{{5.4}{24}{}{ejercicio.5.4}{}}
\newlabel{eje:5-precond@cref}{{[ejercicio][4][5]5.4}{[1][24][]24}}
\newlabel{eje:5-memoria}{{5.5}{24}{}{ejercicio.5.5}{}}
\newlabel{eje:5-memoria@cref}{{[ejercicio][5][5]5.5}{[1][24][]24}}
\newlabel{eje:5-invariancia}{{5.6}{24}{}{ejercicio.5.6}{}}
\newlabel{eje:5-invariancia@cref}{{[ejercicio][6][5]5.6}{[1][24][]24}}
\newlabel{eje:5-sin-correccion}{{5.7}{24}{}{ejercicio.5.7}{}}
\newlabel{eje:5-sin-correccion@cref}{{[ejercicio][7][5]5.7}{[1][24][]24}}
\newlabel{fig:coseno}{{6.1}{27}{Calendario de la \cref {def:coseno} con $\eta _{\max }=3\cdot 10^{-4}$, $\eta _{\min }=3\cdot 10^{-5}$, $W=2000$ y $T=10^{5}$. En el punto medio del tramo coseno, $t=51\,000$, la tasa es $(\eta _{\max }+\eta _{\min })/2=1{,}65\cdot 10^{-4}$}{figure.6.1}{}}
\newlabel{fig:coseno@cref}{{[figure][1][6]6.1}{[1][27][]27}}
\newlabel{ej:coseno}{{6.11}{27}{}{ejemplo.11}{}}
\newlabel{ej:coseno@cref}{{[ejemplo][11][]6.11}{[1][27][]27}}
\newlabel{eje:6-warmup}{{6.1}{28}{}{ejercicio.6.1}{}}
\newlabel{eje:6-warmup@cref}{{[ejercicio][1][6]6.1}{[1][28][]28}}
\newlabel{eje:6-calendario}{{6.2}{28}{}{ejercicio.6.2}{}}
\newlabel{eje:6-calendario@cref}{{[ejercicio][2][6]6.2}{[1][28][]28}}
\newlabel{eje:6-adamw}{{6.3}{28}{}{ejercicio.6.3}{}}
\newlabel{eje:6-adamw@cref}{{[ejercicio][3][6]6.3}{[1][28][]28}}
\newlabel{eje:6-semivida}{{6.4}{28}{}{ejercicio.6.4}{}}
\newlabel{eje:6-semivida@cref}{{[ejercicio][4][6]6.4}{[1][28][]28}}
\newlabel{eje:6-ridge}{{6.5}{28}{}{ejercicio.6.5}{}}
\newlabel{eje:6-ridge@cref}{{[ejercicio][5][6]6.5}{[1][28][]28}}
\newlabel{eje:6-polinomico}{{6.6}{29}{}{ejercicio.6.6}{}}
\newlabel{eje:6-polinomico@cref}{{[ejercicio][6][6]6.6}{[1][28][]29}}
\newlabel{ap:soluciones}{{A}{30}{Soluciones de los ejercicios}{appendix.A}{}}
\newlabel{ap:soluciones@cref}{{[appendix][1][2147483647]A}{[1][30][]30}}
\newlabel{sol:eje:1-cuartica}{{A.1}{30}{Tema 1}{section*.5}{}}
\newlabel{sol:eje:1-cuartica@cref}{{[subappendix][1][2147483647,1]A.1}{[1][30][]30}}
\newlabel{sol:eje:1-silla}{{A.1}{30}{Tema 1}{section*.6}{}}
\newlabel{sol:eje:1-silla@cref}{{[subappendix][1][2147483647,1]A.1}{[1][30][]30}}
\newlabel{sol:eje:1-cuadratica}{{A.1}{31}{Tema 1}{section*.7}{}}
\newlabel{sol:eje:1-cuadratica@cref}{{[subappendix][1][2147483647,1]A.1}{[1][31][]31}}
\newlabel{sol:eje:1-constantes}{{A.1}{31}{Tema 1}{section*.8}{}}
\newlabel{sol:eje:1-constantes@cref}{{[subappendix][1][2147483647,1]A.1}{[1][31][]31}}
\newlabel{sol:eje:1-maximo}{{A.1}{31}{Tema 1}{section*.9}{}}
\newlabel{sol:eje:1-maximo@cref}{{[subappendix][1][2147483647,1]A.1}{[1][31][]31}}
\newlabel{sol:eje:1-fuerte}{{A.1}{32}{Tema 1}{section*.10}{}}
\newlabel{sol:eje:1-fuerte@cref}{{[subappendix][1][2147483647,1]A.1}{[1][32][]32}}
\newlabel{sol:eje:2-escala}{{A.2}{32}{Tema 2}{section*.11}{}}
\newlabel{sol:eje:2-escala@cref}{{[subappendix][2][2147483647,1]A.2}{[1][32][]32}}
\newlabel{sol:eje:2-1d}{{A.2}{33}{Tema 2}{section*.12}{}}
\newlabel{sol:eje:2-1d@cref}{{[subappendix][2][2147483647,1]A.2}{[1][33][]33}}
\newlabel{sol:eje:2-2d}{{A.2}{33}{Tema 2}{section*.13}{}}
\newlabel{sol:eje:2-2d@cref}{{[subappendix][2][2147483647,1]A.2}{[1][33][]33}}
\newlabel{sol:eje:2-valle}{{A.2}{33}{Tema 2}{section*.14}{}}
\newlabel{sol:eje:2-valle@cref}{{[subappendix][2][2147483647,1]A.2}{[1][33][]33}}
\newlabel{sol:eje:2-paso-largo}{{A.2}{34}{Tema 2}{section*.15}{}}
\newlabel{sol:eje:2-paso-largo@cref}{{[subappendix][2][2147483647,1]A.2}{[1][34][]34}}
\newlabel{sol:eje:2-rosenbrock}{{A.2}{34}{Tema 2}{section*.16}{}}
\newlabel{sol:eje:2-rosenbrock@cref}{{[subappendix][2][2147483647,1]A.2}{[1][34][]34}}
\newlabel{sol:eje:2-cuadratica}{{A.2}{35}{Tema 2}{section*.17}{}}
\newlabel{sol:eje:2-cuadratica@cref}{{[subappendix][2][2147483647,1]A.2}{[1][35][]35}}
\newlabel{sol:eje:3-motivos}{{A.3}{36}{Tema 3}{section*.18}{}}
\newlabel{sol:eje:3-motivos@cref}{{[subappendix][3][2147483647,1]A.3}{[1][35][]36}}
\newlabel{sol:eje:3-minibatch}{{A.3}{36}{Tema 3}{section*.19}{}}
\newlabel{sol:eje:3-minibatch@cref}{{[subappendix][3][2147483647,1]A.3}{[1][36][]36}}
\newlabel{sol:eje:3-reemplazo}{{A.3}{36}{Tema 3}{section*.20}{}}
\newlabel{sol:eje:3-reemplazo@cref}{{[subappendix][3][2147483647,1]A.3}{[1][36][]36}}
\newlabel{sol:eje:3-suelo}{{A.3}{37}{Tema 3}{section*.21}{}}
\newlabel{sol:eje:3-suelo@cref}{{[subappendix][3][2147483647,1]A.3}{[1][37][]37}}
\newlabel{sol:eje:3-sin-reemplazo}{{A.3}{37}{Tema 3}{section*.22}{}}
\newlabel{sol:eje:3-sin-reemplazo@cref}{{[subappendix][3][2147483647,1]A.3}{[1][37][]37}}
\newlabel{sol:eje:3-ruido}{{A.3}{38}{Tema 3}{section*.23}{}}
\newlabel{sol:eje:3-ruido@cref}{{[subappendix][3][2147483647,1]A.3}{[1][38][]38}}
\newlabel{sol:eje:4-efectivo}{{A.4}{38}{Tema 4}{section*.24}{}}
\newlabel{sol:eje:4-efectivo@cref}{{[subappendix][4][2147483647,1]A.4}{[1][38][]38}}
\newlabel{sol:eje:4-pasos}{{A.4}{39}{Tema 4}{section*.25}{}}
\newlabel{sol:eje:4-pasos@cref}{{[subappendix][4][2147483647,1]A.4}{[1][39][]39}}
\newlabel{sol:eje:4-optimos}{{A.4}{39}{Tema 4}{section*.26}{}}
\newlabel{sol:eje:4-optimos@cref}{{[subappendix][4][2147483647,1]A.4}{[1][39][]39}}
\newlabel{sol:eje:4-autovalores}{{A.4}{39}{Tema 4}{section*.27}{}}
\newlabel{sol:eje:4-autovalores@cref}{{[subappendix][4][2147483647,1]A.4}{[1][39][]39}}
\newlabel{sol:eje:4-forma}{{A.4}{40}{Tema 4}{section*.28}{}}
\newlabel{sol:eje:4-forma@cref}{{[subappendix][4][2147483647,1]A.4}{[1][40][]40}}
\newlabel{sol:eje:4-estabilidad}{{A.4}{40}{Tema 4}{section*.29}{}}
\newlabel{sol:eje:4-estabilidad@cref}{{[subappendix][4][2147483647,1]A.4}{[1][40][]40}}
\newlabel{sol:eje:5-escala-test}{{A.5}{41}{Tema 5}{section*.30}{}}
\newlabel{sol:eje:5-escala-test@cref}{{[subappendix][5][2147483647,1]A.5}{[1][41][]41}}
\newlabel{sol:eje:5-rmsprop}{{A.5}{41}{Tema 5}{section*.31}{}}
\newlabel{sol:eje:5-rmsprop@cref}{{[subappendix][5][2147483647,1]A.5}{[1][41][]41}}
\newlabel{sol:eje:5-adam}{{A.5}{42}{Tema 5}{section*.32}{}}
\newlabel{sol:eje:5-adam@cref}{{[subappendix][5][2147483647,1]A.5}{[1][42][]42}}
\newlabel{sol:eje:5-precond}{{A.5}{42}{Tema 5}{section*.33}{}}
\newlabel{sol:eje:5-precond@cref}{{[subappendix][5][2147483647,1]A.5}{[1][42][]42}}
\newlabel{sol:eje:5-memoria}{{A.5}{43}{Tema 5}{section*.34}{}}
\newlabel{sol:eje:5-memoria@cref}{{[subappendix][5][2147483647,1]A.5}{[1][43][]43}}
\newlabel{sol:eje:5-invariancia}{{A.5}{43}{Tema 5}{section*.35}{}}
\newlabel{sol:eje:5-invariancia@cref}{{[subappendix][5][2147483647,1]A.5}{[1][43][]43}}
\newlabel{sol:eje:5-sin-correccion}{{A.5}{43}{Tema 5}{section*.36}{}}
\newlabel{sol:eje:5-sin-correccion@cref}{{[subappendix][5][2147483647,1]A.5}{[1][43][]43}}
\newlabel{sol:eje:6-warmup}{{A.6}{44}{Tema 6}{section*.37}{}}
\newlabel{sol:eje:6-warmup@cref}{{[subappendix][6][2147483647,1]A.6}{[1][44][]44}}
\newlabel{sol:eje:6-calendario}{{A.6}{44}{Tema 6}{section*.38}{}}
\newlabel{sol:eje:6-calendario@cref}{{[subappendix][6][2147483647,1]A.6}{[1][44][]44}}
\newlabel{sol:eje:6-adamw}{{A.6}{45}{Tema 6}{section*.39}{}}
\newlabel{sol:eje:6-adamw@cref}{{[subappendix][6][2147483647,1]A.6}{[1][45][]45}}
\newlabel{sol:eje:6-semivida}{{A.6}{45}{Tema 6}{section*.40}{}}
\newlabel{sol:eje:6-semivida@cref}{{[subappendix][6][2147483647,1]A.6}{[1][45][]45}}
\newlabel{sol:eje:6-ridge}{{A.6}{45}{Tema 6}{section*.41}{}}
\newlabel{sol:eje:6-ridge@cref}{{[subappendix][6][2147483647,1]A.6}{[1][45][]45}}
\newlabel{sol:eje:6-polinomico}{{A.6}{46}{Tema 6}{section*.42}{}}
\newlabel{sol:eje:6-polinomico@cref}{{[subappendix][6][2147483647,1]A.6}{[1][46][]46}}
\makeatother

\begin{document}
\frontmatter
\begin{titlepage}
\centering
\vspace*{2.5cm}
{\large\scshape Anatomía de un LLM}\par
\vspace{0.4cm}
{\normalsize Fundamentos matemáticos de los modelos de lenguaje}\par
\vspace{2.2cm}
\rule{\textwidth}{0.8pt}\par\vspace{0.5cm}
{\small\scshape Fase II \textperiodcentered\ Aprendizaje automático y redes neuronales}\par\vspace{0.3cm}
{\Huge\bfseries Módulo 7\par}\vspace{0.3cm}
{\LARGE Optimización\par}
\vspace{0.5cm}\rule{\textwidth}{0.4pt}\par
\vspace{1.6cm}
{\large Resumen del módulo}\par
\vfill
{\small Edición de octubre de 2026}\par
\end{titlepage}

\chapter*{Sobre este resumen}
\addcontentsline{toc}{chapter}{Sobre este resumen}
\markboth{Sobre este resumen}{}

Este documento reúne lo indispensable de los apuntes del Módulo 7: el texto expositivo, todas las definiciones y todos los resultados (teoremas, proposiciones, lemas y corolarios), las tablas y las observaciones de contenido. Se han suprimido las demostraciones, los ejemplos, las lecturas, las figuras, los ejercicios y sus soluciones. La numeración de temas, secciones, enunciados y ecuaciones es la misma que en los apuntes, de modo que cualquier referencia puede buscarse allí; las que remiten a elementos suprimidos (un ejemplo, una figura, una ecuación de una demostración) conservan su número en los apuntes.

\section*{Notación}
\begin{center}
\begin{tabular}{@{}ll@{}}
\toprule
Símbolo & Significado \\
\midrule
$\theta\in\R^{n}$, $\theta_k$ & vector de parámetros e iterado $k$-ésimo \\
$f$, $f^{\ast}=\inf f$, $\theta^{\ast}$ & función objetivo, valor óptimo y minimizador \\
$J(\theta)=\frac1N\sum_{i=1}^{N}\ell_i(\theta)$ & riesgo empírico: pérdida media sobre $N$ ejemplos \\
$\nabla f$, $\nabla^{2}f$ & gradiente (vector columna) y hessiano \\
$\mu$, $L$ & constantes de convexidad fuerte y de suavidad \\
$\lambda_{\min}$, $\lambda_{\max}$, $\lambda_i$ & autovalores de una matriz simétrica \\
$\kappa=L/\mu$ & número de condición \\
$\eta$, $\eta_k$ & tasa de aprendizaje (constante o en el paso $k$) \\
$B$, $\gb$ & tamaño del minibatch y gradiente estocástico \\
$\beta$; $\beta_1$, $\beta_2$ & coeficiente de momentum; coeficientes de Adam \\
$m_t$, $v_t$; $\hat m_t$, $\hat v_t$ & momentos de Adam y sus versiones corregidas \\
$\varepsilon$ & constante de estabilidad numérica de los métodos adaptativos \\
$\lambda$ & coeficiente de regularización o de decaimiento de pesos (tema~6) \\
$A\succeq 0$, $A\succ 0$ & $A$ simétrica semidefinida o definida positiva \\
$u\odot w$, $u^{2}$, $\sqrt{u}$ & operaciones elemento a elemento sobre vectores \\
\bottomrule
\end{tabular}
\end{center}


\tableofcontents

\mainmatter
% =====================================================================
\chapter{Problemas de optimización y convexidad}
\label{tema:1}
% =====================================================================

Este tema plantea el entrenamiento de un modelo como un problema de minimización sin restricciones y reúne las herramientas de análisis que se usan en el resto del módulo: convexidad, convexidad fuerte, suavidad de Lipschitz y condiciones de optimalidad. Las definiciones y los resultados siguen los capítulos~2, 3 y~9 de~\cite{bv}; la discusión sobre redes profundas sigue~\cite[§8.1 y §8.2]{gbc} y~\cite[§12.1 y §12.2]{d2l}. Parte de la materia de convexidad se introdujo en el módulo~4 (tema~5); aquí se completa con las constantes $\mu$ y $L$, que determinan la velocidad de los algoritmos.

\section{Minimización del riesgo empírico}

\begin{definicion}[Problema de optimización sin restricciones]
Sea $f\colon\R^{n}\to\R$. El problema \emph{minimizar $f(\theta)$, $\theta\in\R^{n}$}, consiste en hallar el \emph{valor óptimo} $f^{\ast}=\inf_{\theta}f(\theta)$ y, si existe, un \emph{minimizador global} $\theta^{\ast}$ con $f(\theta^{\ast})=f^{\ast}$. Un punto $\bar\theta$ es un \emph{minimizador local} si existe $r>0$ tal que $f(\bar\theta)\le f(\theta)$ para todo $\theta$ con $\norm{\theta-\bar\theta}<r$.
\end{definicion}

\begin{definicion}[Riesgo empírico]\label{def:riesgo}
Sean $N$ ejemplos de entrenamiento y $\ell_i(\theta)$ la pérdida del modelo de parámetros $\theta$ sobre el ejemplo $i$. El \emph{riesgo empírico} es
\[
J(\theta)=\frac1N\sum_{i=1}^{N}\ell_i(\theta).
\]
\end{definicion}

En un modelo de lenguaje, cada ejemplo es una posición de un texto y $\ell_i(\theta)=-\ln p_\theta(x_t\mid x_{1:t-1})$ es la log-verosimilitud negativa del token observado; $J$ es la pérdida $\mathcal{J}$ del módulo~1 (tema~5). El objetivo real no es minimizar $J$, sino el riesgo esperado sobre la distribución de los datos, al que $J$ solo aproxima; por eso el aprendizaje se diferencia de la optimización pura~\cite[§8.1]{gbc}. Este módulo trata, sin embargo, del problema de optimización.

\section{Convexidad, convexidad fuerte y suavidad}

\begin{definicion}[Conjunto y función convexos]
Un conjunto $C\subseteq\R^{n}$ es \emph{convexo} si $t x+(1-t)y\in C$ para todo $x,y\in C$ y $t\in[0,1]$. Una función $f\colon\R^{n}\to\R$ es \emph{convexa} si
\[
f\bigl(tx+(1-t)y\bigr)\le t f(x)+(1-t)f(y)\qquad\text{para todo } x,y\in\R^{n},\ t\in[0,1],
\]
y \emph{estrictamente convexa} si la desigualdad es estricta cuando $x\ne y$ y $t\in(0,1)$~\cite[cap.~2 y §3.1]{bv}.
\end{definicion}

\begin{teorema}[Caracterización de primer orden]\label{teo:primer-orden}
Sea $f$ diferenciable. Entonces $f$ es convexa si y solo si
\[
f(y)\ge f(x)+\nabla f(x)\T(y-x)\qquad\text{para todo } x,y\in\R^{n}.
\]
\end{teorema}

Geométricamente, una función convexa queda por encima de todas sus rectas (o hiperplanos) tangentes. Para funciones dos veces diferenciables, la convexidad equivale a $\nabla^{2}f(x)\succeq 0$ para todo $x$ (módulo~4, tema~5; \cite[§3.1]{bv}). Las dos definiciones siguientes refuerzan esta condición por abajo y la limitan por arriba.

\begin{definicion}[Convexidad fuerte]\label{def:fuerte}
Sea $\mu>0$. Una función diferenciable $f$ es \emph{$\mu$-fuertemente convexa} si
\[
f(y)\ge f(x)+\nabla f(x)\T(y-x)+\frac{\mu}{2}\norm{y-x}^{2}\qquad\text{para todo } x,y.
\]
Si $f$ es dos veces diferenciable, esto equivale a $\nabla^{2}f(x)\succeq\mu I$ para todo $x$~\cite[§9.1]{bv}.
\end{definicion}

\begin{definicion}[Suavidad de Lipschitz]\label{def:suave}
Sea $L>0$. Una función diferenciable $f$ es \emph{$L$-suave} si su gradiente es lipschitziano de constante $L$:
\[
\norm{\nabla f(x)-\nabla f(y)}\le L\norm{x-y}\qquad\text{para todo } x,y.
\]
Si $f$ es dos veces diferenciable, esto equivale a que todos los autovalores de $\nabla^{2}f(x)$ estén en $[-L,L]$ para todo $x$; para $f$ convexa, a $\nabla^{2}f(x)\preceq LI$.
\end{definicion}

Si $f$ es $\mu$-fuertemente convexa y $L$-suave, necesariamente $\mu\le L$, y el cociente
\[
\kappa=\frac{L}{\mu}\ge 1
\]
se llama \emph{número de condición} de $f$. En el \cref{tema:2} se demuestra que $\kappa$ controla el número de iteraciones del descenso de gradiente.

\addtocounter{definicion}{1}%

\section{Condiciones de optimalidad}

\begin{teorema}[Condición necesaria de primer orden]\label{teo:fermat}
Si $f$ es diferenciable y $\bar\theta$ es un minimizador local, entonces $\nabla f(\bar\theta)=0$.
\end{teorema}

Un punto con $\nabla f(\bar\theta)=0$ se llama \emph{punto crítico} o \emph{estacionario}. En general, la condición no es suficiente; para funciones convexas, sí lo es.

\begin{teorema}[Optimalidad en el caso convexo]\label{teo:opt-convexa}
Sea $f$ convexa y diferenciable. Un punto $\theta^{\ast}$ es minimizador global si y solo si $\nabla f(\theta^{\ast})=0$. En particular, todo minimizador local de una función convexa es global.
\end{teorema}

\begin{proposicion}[Minimizador de una cuadrática]\label{prop:min-cuadratica}
Sea $f(\theta)=\tfrac12\theta\T Q\theta+c\T\theta+b$ con $Q\succ 0$. Su único minimizador y su valor mínimo son
\[
\theta^{\ast}=-Q^{-1}c,\qquad f^{\ast}=b-\tfrac12\,c\T Q^{-1}c,
\]
y para todo $\theta$ se cumple $f(\theta)-f^{\ast}=\tfrac12(\theta-\theta^{\ast})\T Q(\theta-\theta^{\ast})$.
\end{proposicion}

La identidad de la \cref{prop:min-cuadratica} permite trasladar el minimizador al origen: en lo que sigue, al estudiar cuadráticas se supondrá sin pérdida de generalidad $f(\theta)=\tfrac12\theta\T Q\theta$, con $\theta^{\ast}=0$ y $f^{\ast}=0$.

\begin{proposicion}[Desigualdad de Polyak--Łojasiewicz]\label{prop:pl}
Si $f$ es $\mu$-fuertemente convexa y diferenciable, con valor óptimo $f^{\ast}$, entonces
\[
\norm{\nabla f(x)}^{2}\ge 2\mu\bigl(f(x)-f^{\ast}\bigr)\qquad\text{para todo } x.
\]
\end{proposicion}

La \cref{prop:pl} dice que, lejos del óptimo en valor, el gradiente no puede ser pequeño. Es la propiedad que hace converger linealmente al descenso de gradiente (\cref{teo:gd-fuerte}).

\section{Paisajes de pérdida no convexos}

La pérdida de una red neuronal profunda no es convexa en sus parámetros (módulo~4, tema~5). La información local en un punto crítico la proporciona el hessiano.

\begin{proposicion}[Clasificación de puntos críticos]\label{prop:clasificacion}
Sea $f$ de clase $C^{2}$ y $\bar\theta$ un punto crítico.
\begin{enumerate}
  \item Si $\nabla^{2}f(\bar\theta)\succ 0$, $\bar\theta$ es un minimizador local estricto.
  \item Si $\nabla^{2}f(\bar\theta)$ tiene un autovalor negativo, $\bar\theta$ no es minimizador local; si tiene autovalores de ambos signos, es un \emph{punto de silla}.
  \item Si $\nabla^{2}f(\bar\theta)\succeq 0$ con algún autovalor nulo, el criterio no decide.
\end{enumerate}
\end{proposicion}

La demostración, basada en el desarrollo de Taylor de segundo orden, se dio en el módulo~4 (tema~4). El \cref{tab:criticos} recoge los casos con ejemplos en $\R^{2}$.

\begin{table}[h]
\centering\small
\caption{Puntos críticos en el origen y autovalores del hessiano.}\label{tab:criticos}
\begin{tabular}{@{}llll@{}}
\toprule
Tipo & Autovalores de $\nabla^{2}f$ & Ejemplo & Comportamiento del descenso \\
\midrule
Mínimo local & todos $>0$ & $x^{2}+y^{2}$ & converge si el paso es pequeño \\
Máximo local & todos $<0$ & $-x^{2}-y^{2}$ & se aleja \\
Punto de silla & de ambos signos & $x^{2}-y^{2}$ & se aleja salvo en un conjunto de medida nula \\
Valle plano & alguno nulo & $x^{2}$ & no hay fuerza en la dirección plana \\
\bottomrule
\end{tabular}
\end{table}

\addtocounter{definicion}{1}%

\begin{observacion}[El paisaje de las redes profundas]\label{obs:paisaje}
Se conocen varios hechos empíricos sobre la pérdida de las redes profundas. (i) En alta dimensión, los puntos críticos de pérdida alta son mayoritariamente sillas, no mínimos locales: un mínimo exige que todas las curvaturas sean positivas a la vez~\cite{dauphin}, \cite[§8.2.3]{gbc}. (ii) El espectro del hessiano se concentra cerca de cero, con un grupo menor de autovalores alejados que dependen de los datos~\cite{sagun}; el problema está muy mal condicionado~\cite[§8.2.1]{gbc}. (iii) Los lotes muy grandes tienden a llevar a mínimos más agudos, que se han asociado a una peor generalización~\cite{keskar}. (iv) Las conexiones residuales suavizan visiblemente el paisaje~\cite{li-paisaje}; se estudian en el módulo~11. Los temas siguientes desarrollan algoritmos que funcionan razonablemente en este entorno aunque su análisis se haga en el caso convexo.
\end{observacion}


% =====================================================================
\chapter{Descenso de gradiente}
\label{tema:2}
% =====================================================================

El descenso de gradiente es el algoritmo del que derivan todos los optimizadores del módulo. Este tema demuestra sus garantías de convergencia en tres situaciones de generalidad decreciente: funciones suaves arbitrarias, funciones fuertemente convexas y funciones cuadráticas, en las que el análisis es exacto. El resultado central es que el número de iteraciones crece proporcionalmente al número de condición $\kappa$. Se sigue~\cite[§9.2 y §9.3]{bv}, \cite[§12.3]{d2l} y las clases~21 y~22 de~\cite{mit18065}.

\section{El método y el lema de descenso}

\begin{definicion}[Descenso de gradiente]\label{def:gd}
Dados $\theta_0\in\R^{n}$ y una \emph{tasa de aprendizaje} (o \emph{paso}) $\eta>0$, el \emph{descenso de gradiente} genera la sucesión
\[
\theta_{k+1}=\theta_k-\eta\,\nabla f(\theta_k),\qquad k=0,1,2,\dots
\]
\end{definicion}

La dirección $-\nabla f(\theta_k)$ es la de máximo descenso local (módulo~4, tema~1). Por el desarrollo de Taylor de primer orden, $f(\theta_k-\eta\nabla f(\theta_k))=f(\theta_k)-\eta\norm{\nabla f(\theta_k)}^{2}+o(\eta)$, de modo que un paso suficientemente pequeño reduce $f$ siempre que el gradiente no sea nulo. El lema siguiente cuantifica «suficientemente pequeño».

\begin{lema}[Lema de descenso]\label{lem:descenso}
Si $f$ es $L$-suave, para todo $x,y\in\R^{n}$,
\[
f(y)\le f(x)+\nabla f(x)\T(y-x)+\frac{L}{2}\norm{y-x}^{2}.
\]
\end{lema}

\begin{corolario}[Descenso suficiente]\label{cor:descenso}
Si $f$ es $L$-suave y $\theta^{+}=\theta-\eta\nabla f(\theta)$ con $0<\eta<2/L$, entonces
\[
f(\theta^{+})\le f(\theta)-\eta\Bigl(1-\frac{L\eta}{2}\Bigr)\norm{\nabla f(\theta)}^{2}.
\]
Con $\eta=1/L$ la cota es $f(\theta^{+})\le f(\theta)-\frac{1}{2L}\norm{\nabla f(\theta)}^{2}$, y es la mejor que da el lema.
\end{corolario}

\begin{teorema}[Caso no convexo]\label{teo:gd-noconvexo}
Sea $f$ $L$-suave y acotada inferiormente por $f^{\ast}$. El descenso de gradiente con $\eta=1/L$ cumple, para todo $K\ge 1$,
\[
\min_{0\le k<K}\norm{\nabla f(\theta_k)}^{2}\le\frac{2L\bigl(f(\theta_0)-f^{\ast}\bigr)}{K}.
\]
\end{teorema}

El \cref{teo:gd-noconvexo} solo garantiza aproximarse a un punto crítico, que puede ser una silla; es la clase de garantía que cabe esperar para la pérdida de una red profunda.

\section{Convergencia lineal en funciones fuertemente convexas}

\begin{teorema}[Convergencia lineal]\label{teo:gd-fuerte}
Sea $f$ $\mu$-fuertemente convexa y $L$-suave, con valor óptimo $f^{\ast}$. El descenso de gradiente con $\eta=1/L$ cumple
\[
f(\theta_k)-f^{\ast}\le\Bigl(1-\frac{1}{\kappa}\Bigr)^{k}\bigl(f(\theta_0)-f^{\ast}\bigr),\qquad\kappa=\frac{L}{\mu}.
\]
\end{teorema}

\begin{corolario}[Número de iteraciones]\label{cor:iteraciones}
En las condiciones del \cref{teo:gd-fuerte}, para alcanzar $f(\theta_k)-f^{\ast}\le\epsilon\,(f(\theta_0)-f^{\ast})$ basta con $k\ge\kappa\ln(1/\epsilon)$ iteraciones.
\end{corolario}

La convergencia se llama \emph{lineal} porque el error se multiplica en cada paso por un factor constante menor que~1: el logaritmo del error decrece linealmente con $k$. El factor depende solo de $\kappa$. El análisis de~\cite[§9.3]{bv} obtiene el mismo factor $1-\mu/L$, salvo constantes, para búsquedas lineales exacta y con retroceso.

\section{Análisis exacto en cuadráticas}

En una función cuadrática el método es lineal en $\theta$ y puede resolverse en forma cerrada. Toda función de clase $C^{2}$ es aproximadamente cuadrática cerca de un mínimo no degenerado, de modo que este caso describe el comportamiento asintótico general.

\addtocounter{definicion}{1}%

\begin{teorema}[Descenso de gradiente en cuadráticas]\label{teo:gd-cuadratica}
Sea $f(\theta)=\tfrac12\theta\T Q\theta$ con $Q$ simétrica de autovalores $0<\mu=\lambda_1\le\dots\le\lambda_n=L$. Sea $\rho(\eta)=\max_i\abs{1-\eta\lambda_i}$.
\begin{enumerate}
  \item $\norm{\theta_k}\le\rho(\eta)^{k}\norm{\theta_0}$, con igualdad si $\theta_0$ es autovector del autovalor que realiza el máximo.
  \item El método converge a $\theta^{\ast}=0$ para todo $\theta_0$ si y solo si $0<\eta<2/L$.
  \item El valor de $\eta$ que minimiza $\rho(\eta)$ es $\eta^{\ast}=2/(\mu+L)$, y
  \[
  \rho(\eta^{\ast})=\frac{L-\mu}{L+\mu}=\frac{\kappa-1}{\kappa+1}.
  \]
\end{enumerate}
\end{teorema}

\begin{observacion}
El \cref{teo:gd-cuadratica} explica la conducta del método. La dirección de curvatura máxima, $\lambda_n=L$, limita el paso: $\eta<2/L$. La dirección de curvatura mínima, $\lambda_1=\mu$, se contrae entonces con un factor $1-\eta\mu>1-2/\kappa$, próximo a~1 si $\kappa$ es grande. Con el paso óptimo, las dos direcciones extremas se contraen al mismo ritmo, pero con signos opuestos: la de curvatura alta oscila.
\end{observacion}

\addtocounter{definicion}{1}%

\addtocounter{figure}{1}%

\begin{corolario}\label{cor:iter-cuadratica}
En las condiciones del \cref{teo:gd-cuadratica}, con el paso óptimo, el menor $k$ con $\rho^{k}\le\epsilon$ es $k=\lceil\ln\epsilon/\ln\rho\rceil$, que para $\kappa$ grande es aproximadamente $\frac{\kappa}{2}\ln(1/\epsilon)$.
\end{corolario}

El \cref{tab:kappa} aplica el corolario con $\epsilon=10^{-3}$. La última columna anticipa el resultado del \cref{tema:4}: con momentum, el número de iteraciones crece como $\sqrt\kappa$ en lugar de como $\kappa$.

\begin{table}[h]
\centering\small
\caption{Factor de contracción e iteraciones necesarias para reducir el error a $10^{-3}$ en una cuadrática, con parámetros óptimos.}\label{tab:kappa}
\begin{tabular}{@{}rcrcr@{}}
\toprule
$\kappa$ & $\rho$ del gradiente & Iteraciones & $\rho$ con momentum (\cref{teo:heavy-ball}) & Iteraciones \\
\midrule
10 & $0{,}8182$ & 35 & $0{,}5195$ & 11 \\
100 & $0{,}9802$ & 346 & $0{,}8182$ & 35 \\
1000 & $0{,}9980$ & 3454 & $0{,}9387$ & 110 \\
\bottomrule
\end{tabular}
\end{table}

\section{La función de Rosenbrock y el borde de estabilidad}

\addtocounter{definicion}{1}%

\begin{observacion}[Borde de estabilidad]\label{obs:borde}
La condición $\eta<2/\lambda_{\max}$ también se observa en redes profundas, pero en sentido inverso. Cohen \emph{et al.}~\cite{cohen} observaron que, al entrenar con descenso de gradiente de paso fijo, la curvatura máxima $\lambda_{\max}$ del hessiano de la pérdida crece hasta situarse en torno a $2/\eta$ y permanece ahí; la pérdida deja de decrecer de forma monótona a corto plazo, pero sigue bajando a largo plazo. La red evoluciona hacia el límite de estabilidad del \cref{teo:gd-cuadratica}.
\end{observacion}


% =====================================================================
\chapter{Descenso de gradiente estocástico}
\label{tema:3}
% =====================================================================

Un modelo de lenguaje se entrena con billones de tokens; calcular el gradiente exacto del riesgo empírico antes de cada paso es inviable. El descenso de gradiente estocástico sustituye ese gradiente por una estimación barata obtenida de un subconjunto aleatorio de ejemplos. Este tema estudia las propiedades estadísticas de esa estimación, demuestra una garantía de convergencia en el caso convexo y explica por qué el ruido obliga a reducir la tasa de aprendizaje al final del entrenamiento. Se sigue~\cite[§12.4 y §12.5]{d2l}, \cite[§8.1.3 y §8.3.1]{gbc}, la revisión~\cite{bottou} y la clase~25 de~\cite{mit18065}.

\section{El gradiente de minibatch}

Sea $J(\theta)=\frac1N\sum_{i=1}^{N}\ell_i(\theta)$ el riesgo empírico (\cref{def:riesgo}), con cada $\ell_i$ diferenciable.

\begin{definicion}[Descenso de gradiente estocástico]\label{def:sgd}
Sea $B\ge 1$. En cada paso $k$ se eligen índices $I_1,\dots,I_B$ en $\{1,\dots,N\}$, independientes entre sí y de los pasos anteriores, y se calcula el \emph{gradiente de minibatch}
\[
\gb_k=\frac1B\sum_{b=1}^{B}\nabla\ell_{I_b}(\theta_k).
\]
El \emph{descenso de gradiente estocástico} (SGD) es la iteración $\theta_{k+1}=\theta_k-\eta_k\,\gb_k$. Con $B=1$ se obtiene el SGD en sentido estricto; con $B=N$ y sin aleatoriedad, el descenso de gradiente del \cref{tema:2}.
\end{definicion}

Para un vector aleatorio $X\in\R^{n}$ se usa la \emph{varianza total} $\Var[X]=\E\norm{X-\E X}^{2}$, suma de las varianzas de las componentes. Sea, para $\theta$ fijo,
\[
\sigma^{2}(\theta)=\frac1N\sum_{i=1}^{N}\norm{\nabla\ell_i(\theta)-\nabla J(\theta)}^{2},
\]
la varianza de los gradientes individuales.

\begin{teorema}[Insesgadez y varianza del minibatch]\label{teo:minibatch}
Si $I_1,\dots,I_B$ son independientes y uniformes en $\{1,\dots,N\}$ (muestreo con reemplazamiento), entonces, para $\theta$ fijo,
\[
\E[\gb]=\nabla J(\theta),\qquad\Var[\gb]=\frac{\sigma^{2}(\theta)}{B}.
\]
\end{teorema}

En la práctica se recorre el conjunto de datos en orden aleatorio, lo que equivale a muestrear sin reemplazamiento dentro de cada época. La varianza es entonces algo menor.

\begin{proposicion}[Muestreo sin reemplazamiento]\label{prop:sin-reemplazo}
Si $\{I_1,\dots,I_B\}$ es un subconjunto de $B$ índices distintos elegido uniformemente entre los $\binom{N}{B}$ posibles, con $1\le B\le N$, entonces $\E[\gb]=\nabla J(\theta)$ y
\[
\Var[\gb]=\frac{\sigma^{2}(\theta)}{B}\cdot\frac{N-B}{N-1}.
\]
\end{proposicion}

La demostración se propone en el \cref{eje:3-sin-reemplazo}. El factor $(N-B)/(N-1)$ es la \emph{corrección por población finita} de la estadística clásica: vale~1 si $B=1$ y~0 si $B=N$, caso en que $\gb$ es el gradiente exacto.

\addtocounter{definicion}{1}%

\begin{observacion}[Coste y tamaño del lote]
La desviación típica del gradiente decrece como $1/\sqrt{B}$, mientras que el coste de calcularlo crece como $B$: duplicar el lote solo divide el ruido por $\sqrt2\approx 1{,}414$. Por eso suele compensar dar muchos pasos ruidosos y baratos antes que pocos exactos~\cite[§8.1.3]{gbc}. El tamaño del lote está limitado además por el paralelismo del hardware (módulo~17). Los modelos de lenguaje usan lotes del orden de millones de tokens: GPT-3, en su versión de $175\cdot 10^{9}$ parámetros, se entrenó con lotes de $3{,}2$ millones de tokens~\cite[tabla~2.1]{gpt3}. Al aumentar el lote, Goyal \emph{et al.}~\cite{goyal} propusieron escalar la tasa de aprendizaje en la misma proporción (\emph{regla de escalado lineal}), acompañada de un calentamiento inicial.
\end{observacion}

\addtocounter{figure}{1}%

\section{Convergencia en el caso convexo}

\begin{teorema}[SGD en funciones convexas]\label{teo:sgd-convexo}
Sea $f$ convexa y diferenciable, con minimizador $\theta^{\ast}$. Sea la iteración $\theta_{t+1}=\theta_t-\eta_t g_t$, $t=1,\dots,T$, donde $g_t$ es un gradiente estocástico que, condicionado a $\theta_t$, cumple $\E[g_t\mid\theta_t]=\nabla f(\theta_t)$ y $\E[\norm{g_t}^{2}\mid\theta_t]\le G^{2}$. Sean $r=\norm{\theta_1-\theta^{\ast}}$ y $\bar\theta_T=\sum_t\eta_t\theta_t/\sum_t\eta_t$. Entonces
\[
\E\bigl[f(\bar\theta_T)\bigr]-f(\theta^{\ast})\le\frac{r^{2}+G^{2}\sum_{t=1}^{T}\eta_t^{2}}{2\sum_{t=1}^{T}\eta_t}.
\]
\end{teorema}

\begin{corolario}\label{cor:sgd-constante}
Con $T$ fijado y paso constante $\eta_t=\eta=r/(G\sqrt T)$, la cota del \cref{teo:sgd-convexo} vale $rG/\sqrt T$.
\end{corolario}

\begin{corolario}[Condiciones de Robbins--Monro]\label{cor:robbins}
Si $\sum_{t\ge 1}\eta_t=\infty$ y $\sum_{t\ge 1}\eta_t^{2}<\infty$, la cota del \cref{teo:sgd-convexo} tiende a~0 cuando $T\to\infty$.
\end{corolario}

Las dos condiciones del \cref{cor:robbins} proceden del trabajo fundacional de Robbins y Monro sobre aproximación estocástica~\cite{robbins}: los pasos deben sumar lo suficiente para llegar a cualquier punto, pero sus cuadrados, que miden el ruido acumulado, deben ser sumables. La velocidad $O(1/\sqrt T)$ es mucho más lenta que la convergencia lineal del \cref{teo:gd-fuerte}: es el precio del ruido.

\section{El ruido y la tasa de aprendizaje}

Con un paso constante, el SGD no converge al minimizador, sino que fluctúa a su alrededor. La proposición siguiente lo cuantifica en el caso más sencillo.

\begin{proposicion}[Suelo de ruido]\label{prop:suelo}
Sea $f(x)=\frac{a}{2}x^{2}$ con $a>0$ y la iteración $x_{k+1}=x_k-\eta(ax_k+\xi_k)$, donde las $\xi_k$ son variables aleatorias de media~0 y varianza $s^{2}$, independientes entre sí y de $x_0$. Si $0<\eta a<2$, entonces $\E[x_k]\to 0$ y
\[
\lim_{k\to\infty}\E[x_k^{2}]=\frac{\eta\,s^{2}}{a(2-\eta a)},\qquad
\lim_{k\to\infty}\E[f(x_k)]=\frac{\eta\,s^{2}}{2(2-\eta a)}.
\]
\end{proposicion}

El exceso de pérdida estacionario es aproximadamente $\eta s^{2}/4$ si $\eta a\ll 1$: proporcional a la tasa de aprendizaje y a la varianza del gradiente. Reducir $\eta$ al final del entrenamiento rebaja ese suelo; es la justificación del decaimiento de la tasa de aprendizaje que se estudia en el \cref{tema:6}.

\begin{observacion}[El ruido no es solo un inconveniente]
En problemas no convexos, el ruido del minibatch tiene efectos favorables: una iteración exactamente sobre la variedad estable de una silla, como en el \cref{ej:silla}, se desvía de ella con probabilidad~1, y los lotes pequeños tienden a encontrar mínimos más anchos~\cite{keskar}. Estos efectos se conocen sobre todo empíricamente; su análisis riguroso es materia de investigación.
\end{observacion}


% =====================================================================
\chapter{Momentum}
\label{tema:4}
% =====================================================================

El \cref{tema:2} mostró que el descenso de gradiente es lento en valles alargados: la dirección de curvatura alta obliga a usar pasos pequeños y la de curvatura baja avanza despacio. El método de momentum, propuesto por Polyak~\cite{polyak} con el nombre de método de la bola pesada, corrige este defecto acumulando los gradientes pasados. Este tema lo presenta como una media móvil exponencial, lo interpreta como la discretización de un sistema mecánico con rozamiento y demuestra que, en cuadráticas, reduce el número de iteraciones de un orden $\kappa$ a un orden $\sqrt\kappa$. Se sigue~\cite[§12.6]{d2l}, \cite[§8.3.2 y §8.3.3]{gbc}, la clase~23 de~\cite{mit18065} y el artículo divulgativo de Goh~\cite{goh}.

\section{Medias móviles exponenciales}

\begin{definicion}[Media móvil exponencial]\label{def:ema}
Sean $g_1,g_2,\dots$ vectores y $\beta\in[0,1)$. La \emph{acumulación exponencial} con $v_0=0$ es $v_k=\beta v_{k-1}+g_k$, y la \emph{media móvil exponencial} es $\bar v_k=\beta\bar v_{k-1}+(1-\beta)g_k$, con $\bar v_0=0$.
\end{definicion}

\begin{proposicion}\label{prop:ema}
Con la notación de la \cref{def:ema},
\[
v_k=\sum_{j=1}^{k}\beta^{k-j}g_j,\qquad \bar v_k=(1-\beta)\,v_k.
\]
En particular, si $g_j=g$ para todo $j$, entonces $v_k=\frac{1-\beta^{k}}{1-\beta}\,g\to\frac{g}{1-\beta}$ y $\bar v_k=(1-\beta^{k})\,g\to g$.
\end{proposicion}

El peso de un gradiente decae geométricamente con su antigüedad; la suma total de los pesos es $1/(1-\beta)$, que se interpreta como el número efectivo de gradientes promediados~\cite[§12.6.1]{d2l}. Con $\beta=0{,}9$ son unos~10; con $\beta=0{,}99$, unos~100.

\section{El método de la bola pesada}

\begin{definicion}[Momentum]\label{def:momentum}
Dados $\theta_0$, $v_0=0$, un paso $\eta>0$ y un coeficiente $\beta\in[0,1)$, el método de \emph{momentum} (o de la \emph{bola pesada}) es
\[
v_{k+1}=\beta v_k+\nabla f(\theta_k),\qquad \theta_{k+1}=\theta_k-\eta\,v_{k+1}.
\]
Con $\beta=0$ se recupera el descenso de gradiente.
\end{definicion}

\begin{proposicion}[Forma de dos pasos]\label{prop:dos-pasos}
Los iterados del método de momentum cumplen, para $k\ge 1$,
\[
\theta_{k+1}=\theta_k-\eta\nabla f(\theta_k)+\beta\,(\theta_k-\theta_{k-1}).
\]
\end{proposicion}

El término $\beta(\theta_k-\theta_{k-1})$ es la inercia: el iterado conserva una fracción $\beta$ de su desplazamiento anterior. Si el gradiente es constante, la \cref{prop:ema} da $\eta v_k\to\eta g/(1-\beta)$: con $\beta=0{,}9$, el paso efectivo es diez veces el del descenso de gradiente. En un valle estrecho, las componentes del gradiente que cambian de signo de un paso al siguiente (las transversales al valle) se cancelan parcialmente en $v_k$, mientras que las que conservan el signo (a lo largo del valle) se acumulan.

\addtocounter{definicion}{1}%

\section{Interpretación física}

Considérese una partícula de masa $m>0$ que se mueve en un potencial $f$ con rozamiento viscoso de coeficiente $\gamma>0$. Su posición $\theta(t)$ y su velocidad $u(t)=\dot\theta(t)$ cumplen
\begin{equation}\label{eq:bola}
m\,\dot u=-\nabla f(\theta)-\gamma\,u,\qquad\dot\theta=u.
\end{equation}

\begin{proposicion}[Disipación de energía]\label{prop:energia}
Sea $E(t)=\frac{m}{2}\norm{u(t)}^{2}+f(\theta(t))$ la energía mecánica. A lo largo de las soluciones de~\eqref{eq:bola}, $\dot E=-\gamma\norm{u}^{2}\le 0$.
\end{proposicion}

La energía decrece mientras la partícula se mueve y solo se estabiliza cuando $u=0$, es decir, en reposo en un punto crítico. Sin rozamiento, la energía se conserva y la partícula oscila indefinidamente.

\begin{proposicion}[Discretización]\label{prop:discretizacion}
Se discretiza~\eqref{eq:bola} con paso temporal $\Delta t>0$ mediante
\[
u_{k+1}=u_k+\frac{\Delta t}{m}\bigl(-\nabla f(\theta_k)-\gamma u_k\bigr),\qquad\theta_{k+1}=\theta_k+\Delta t\,u_{k+1}.
\]
Con $v_k=-(m/\Delta t)\,u_k$, la iteración es exactamente el método de momentum de la \cref{def:momentum} con
\[
\beta=1-\frac{\gamma\,\Delta t}{m},\qquad\eta=\frac{\Delta t^{2}}{m}.
\]
\end{proposicion}

El coeficiente $\beta$ mide la fracción de velocidad que sobrevive a un paso: $\beta\to 1$ corresponde a rozamiento nulo y $\beta=0$ a un rozamiento tan grande que la partícula olvida su velocidad en cada paso, que es el descenso de gradiente.

\section{Convergencia en cuadráticas}

\begin{lema}[Una dirección propia]\label{lem:polinomio}
Sea $f(z)=\frac{\lambda}{2}z^{2}$ con $\lambda>0$. Los iterados del método de momentum cumplen $z_{k+1}=(1+\beta-\eta\lambda)\,z_k-\beta z_{k-1}$, cuyo polinomio característico es
\[
p(r)=r^{2}-(1+\beta-\eta\lambda)\,r+\beta.
\]
Si $(1-\sqrt\beta)^{2}\le\eta\lambda\le(1+\sqrt\beta)^{2}$, las dos raíces de $p$ tienen módulo $\sqrt\beta$.
\end{lema}

\begin{teorema}[Momentum óptimo en cuadráticas]\label{teo:heavy-ball}
Sea $f(\theta)=\frac12\theta\T Q\theta$ con $Q$ simétrica de autovalores en $[\mu,L]$, $0<\mu<L$, y $\kappa=L/\mu$. Con
\[
\eta=\frac{4}{(\sqrt L+\sqrt\mu)^{2}},\qquad\beta=\Bigl(\frac{\sqrt\kappa-1}{\sqrt\kappa+1}\Bigr)^{2},
\]
la iteración lineal del método de momentum sobre el par $(\theta_k,\theta_{k-1})$ tiene radio espectral
\[
\rho=\sqrt\beta=\frac{\sqrt\kappa-1}{\sqrt\kappa+1}.
\]
\end{teorema}

\begin{observacion}
El radio espectral determina la tasa asintótica: para todo $\delta>0$ existe $C$ tal que $\norm{\theta_k}\le C(\rho+\delta)^{k}$. A diferencia del \cref{teo:gd-cuadratica}, la matriz de la iteración no es simétrica y la constante $C$ puede ser mayor que~1; en los extremos $\lambda=\mu$ y $\lambda=L$ la raíz es doble y aparece un factor polinómico en $k$. Por eso las iteraciones reales pueden ser algo más que las que predice $\rho$ (\cref{eje:4-autovalores}).
\end{observacion}

Comparando con el \cref{teo:gd-cuadratica}, para $\kappa$ grande, $\ln\rho_{\text{mom}}\approx -2/\sqrt\kappa$ frente a $\ln\rho_{\text{gd}}\approx -2/\kappa$: el número de iteraciones pasa de ser proporcional a $\kappa$ a serlo a $\sqrt\kappa$. Con $\kappa=100$, el factor baja de $0{,}9802$ a $0{,}8182$ y las iteraciones para reducir el error a $10^{-3}$, de 346 a 35 (\cref{tab:kappa}). El artículo de Goh~\cite{goh} permite explorar este resultado de forma interactiva.

\addtocounter{definicion}{1}%

\section{Momentum de Nesterov}

\begin{definicion}[Momentum de Nesterov]
El método de \emph{momentum de Nesterov} evalúa el gradiente en el punto al que llevaría la inercia:
\[
v_{k+1}=\beta v_k+\nabla f(\theta_k-\eta\beta v_k),\qquad\theta_{k+1}=\theta_k-\eta v_{k+1}.
\]
\end{definicion}

Para funciones convexas y $L$-suaves, una versión de este método con $\beta$ variable alcanza $f(\theta_k)-f^{\ast}=O(1/k^{2})$, frente al $O(1/k)$ del descenso de gradiente, y ese orden es óptimo entre los métodos de primer orden~\cite{nesterov}. En el caso estocástico, en cambio, el momentum de Nesterov no mejora el orden de convergencia~\cite[§8.3.3]{gbc}. En el entrenamiento de redes profundas, el momentum se usa sobre todo a través de Adam (\cref{tema:5}), cuyo primer momento es una media móvil exponencial del gradiente.


% =====================================================================
\chapter{Métodos adaptativos: AdaGrad, RMSProp y Adam}
\label{tema:5}
% =====================================================================

El momentum acelera el descenso en valles mal condicionados, pero usa la misma tasa de aprendizaje en todas las coordenadas. Los métodos adaptativos asignan a cada parámetro su propio paso, dividiendo su gradiente por una estimación de su escala típica. Este tema los presenta como preacondicionadores diagonales, estudia AdaGrad y RMSProp y analiza con detalle Adam, el optimizador sobre el que se construye la receta de preentrenamiento de los modelos de lenguaje. Se sigue~\cite[§12.7, §12.8 y §12.10]{d2l}, \cite[§8.5]{gbc} y los artículos originales~\cite{adagrad,rmsprop,adam}.

En todo el tema, las operaciones entre vectores ($g^{2}$, $\sqrt{s}$, cocientes) se entienden elemento a elemento.

\section{Preacondicionamiento diagonal}

\begin{definicion}[Descenso preacondicionado]\label{def:precond}
Sea $D\in\R^{n\times n}$ diagonal con entradas positivas. El \emph{descenso de gradiente preacondicionado} por $D$ es $\theta_{k+1}=\theta_k-\eta\,D^{-1}\nabla f(\theta_k)$.
\end{definicion}

\begin{proposicion}\label{prop:precond}
El descenso preacondicionado por $D$ sobre $f$ es el descenso de gradiente ordinario sobre $\tilde f(\varphi)=f(D^{-1/2}\varphi)$ en las variables $\varphi=D^{1/2}\theta$. Si $f(\theta)=\frac12\theta\T Q\theta$, el hessiano de $\tilde f$ es $D^{-1/2}QD^{-1/2}$.
\end{proposicion}

Si $Q=\diag(a_1,\dots,a_n)$ y se toma $D=Q$, el hessiano transformado es la identidad y $\kappa=1$: el método converge en un paso con $\eta=1$. Un buen preacondicionador diagonal iguala las curvaturas de las coordenadas. Pero solo puede corregir el mal condicionamiento alineado con los ejes; si los autovectores de $Q$ están girados respecto de ellos, una matriz diagonal no basta (\cref{eje:5-precond}). Los métodos de este tema estiman $D$ a partir de los propios gradientes, sin calcular el hessiano.

\section{AdaGrad y RMSProp}

\begin{definicion}[AdaGrad]\label{def:adagrad}
Con $s_0=0$ y $\varepsilon>0$ pequeño, \emph{AdaGrad}~\cite{adagrad} acumula los cuadrados de los gradientes:
\[
s_t=s_{t-1}+g_t^{2},\qquad\theta_t=\theta_{t-1}-\eta\,\frac{g_t}{\sqrt{s_t}+\varepsilon}.
\]
\end{definicion}

\begin{proposicion}\label{prop:adagrad}
Si $\varepsilon=0$ y una coordenada recibe un gradiente constante $g\ne 0$, su paso en la iteración $t$ mide $\eta/\sqrt t$.
\end{proposicion}

El paso de AdaGrad no depende de la escala del gradiente y decrece como $1/\sqrt t$, un calendario incorporado que es adecuado en problemas convexos pero demasiado agresivo en el entrenamiento de redes profundas, donde la acumulación sin olvido acaba paralizando el método~\cite[§8.5.1]{gbc}.

\begin{definicion}[RMSProp]\label{def:rmsprop}
\emph{RMSProp}~\cite{rmsprop} sustituye la suma por una media móvil exponencial de coeficiente $\rho\in(0,1)$:
\[
s_t=\rho\,s_{t-1}+(1-\rho)\,g_t^{2},\qquad\theta_t=\theta_{t-1}-\eta\,\frac{g_t}{\sqrt{s_t}+\varepsilon}.
\]
\end{definicion}

$\sqrt{s_t}$ estima la raíz de la media cuadrática reciente del gradiente de cada coordenada. Una coordenada con gradientes grandes recibe pasos pequeños, y viceversa. Por la \cref{prop:ema}, con gradiente constante $s_t=(1-\rho^{t})g^{2}$, de modo que el primer paso mide $\eta/\sqrt{1-\rho}$: unas $3{,}16$ veces $\eta$ con $\rho=0{,}9$. Este sesgo inicial, producido por el arranque en $s_0=0$, es lo que corrige Adam.

\section{Adam}

\begin{definicion}[Adam]\label{def:adam}
Sean $\eta>0$, $\beta_1,\beta_2\in[0,1)$, $\varepsilon>0$ y $m_0=v_0=0$. Para $t=1,2,\dots$, con $g_t$ el gradiente (estocástico) en $\theta_{t-1}$, \emph{Adam}~\cite[algoritmo~1]{adam} realiza
\begin{align*}
m_t&=\beta_1m_{t-1}+(1-\beta_1)\,g_t, & \hat m_t&=\frac{m_t}{1-\beta_1^{t}},\\
v_t&=\beta_2v_{t-1}+(1-\beta_2)\,g_t^{2}, & \hat v_t&=\frac{v_t}{1-\beta_2^{t}},\\
\theta_t&=\theta_{t-1}-\eta\,\frac{\hat m_t}{\sqrt{\hat v_t}+\varepsilon}.
\end{align*}
Los valores recomendados en el artículo original son $\eta=10^{-3}$, $\beta_1=0{,}9$, $\beta_2=0{,}999$ y $\varepsilon=10^{-8}$.
\end{definicion}

Adam combina el momentum (primer momento $m_t$, una media móvil del gradiente) con la normalización de RMSProp (segundo momento $v_t$, una media móvil de su cuadrado). Las divisiones por $1-\beta_1^{t}$ y $1-\beta_2^{t}$ son la \emph{corrección de sesgo}~\cite[§3]{adam}.

\begin{proposicion}[Corrección de sesgo]\label{prop:sesgo}
Si $g_t=g$ para todo $t$, entonces $m_t=(1-\beta_1^{t})\,g$ y $\hat m_t=g$. Más en general, si los $g_t$ son aleatorios con $\E[g_t]=\bar g$ para todo $t$, entonces $\E[m_t]=(1-\beta_1^{t})\,\bar g$ y $\E[\hat m_t]=\bar g$. Los mismos resultados valen para $v_t$ y $\hat v_t$ con $g_t^{2}$ y $\beta_2$.
\end{proposicion}

Sin corrección, $m_t$ y $v_t$ estarían sesgados hacia~0 al principio. Con $\beta_1=0{,}9$, $1-\beta_1^{t}$ vale $0{,}1$ en $t=1$, $0{,}19$ en $t=2$ y $0{,}6513$ en $t=10$; con $\beta_2=0{,}999$, $1-\beta_2^{t}$ vale $0{,}001$ en $t=1$ y solo $0{,}6323$ en $t=1000$: el sesgo de $v_t$ persiste durante unos miles de pasos.

\begin{proposicion}[Primer paso e invariancia de escala]\label{prop:escala}
Sea $\varepsilon=0$. (i) En $t=1$, cada coordenada con $g_1\ne 0$ se desplaza exactamente $\eta\,\sign(g_1)$. (ii) Si todos los gradientes se multiplican por una constante $c>0$, los iterados de Adam no cambian.
\end{proposicion}

La invariancia de escala significa que $\eta$ fija directamente el tamaño del paso en el espacio de parámetros: Kingma y Ba observan que el desplazamiento de cada coordenada queda aproximadamente acotado por $\eta$~\cite[§2.1]{adam}. Esto hace a Adam mucho menos sensible que el SGD a la escala de la pérdida y a las diferencias de escala entre capas. Con $\varepsilon>0$ la invariancia es solo aproximada; $\varepsilon$ evita divisiones por cero cuando $\hat v_t\approx 0$.

\addtocounter{definicion}{1}%

\addtocounter{figure}{1}%

\begin{observacion}[Adam en los modelos de lenguaje]
Los artículos de modelos de lenguaje abiertos documentan variantes de Adam con $\beta_2$ menor que el valor recomendado: LLaMA usa AdamW con $\beta_1=0{,}9$ y $\beta_2=0{,}95$~\cite[§2.3]{llama}. Con $\beta_2=0{,}95$, el segundo momento promedia efectivamente unos $1/(1-\beta_2)=20$ pasos en lugar de~1000, y se adapta antes a cambios bruscos en la escala del gradiente. El precio de Adam es la memoria: guarda $m$ y $v$, dos números adicionales por parámetro. En un modelo de $7\cdot 10^{9}$ parámetros son $1{,}4\cdot 10^{10}$ valores, 56~GB en precisión simple (módulo~17). Los optimizadores usados en Claude, Gemini y Grok no se han publicado.
\end{observacion}


% =====================================================================
\chapter{Regularización y calendarios de la tasa de aprendizaje}
\label{tema:6}
% =====================================================================

Este tema completa la receta de optimización de los modelos de lenguaje actuales con dos ingredientes que no cambian la dirección del paso, sino su forma y su tamaño: el decaimiento de pesos y el calendario de la tasa de aprendizaje. Se estudia primero la regularización $L^{2}$ y su relación con el decaimiento de pesos, que es una equivalencia para el SGD pero no para Adam; de ahí surge AdamW. Después se definen el calentamiento lineal y el decaimiento coseno y se justifican con los resultados de los temas~3 y~5. Se sigue~\cite[§7.1]{gbc}, \cite[§12.11]{d2l} y los artículos~\cite{adamw,sgdr,goyal}.

\section{Regularización \texorpdfstring{$L^{2}$}{L2}}

\begin{definicion}[Regularización $L^{2}$]\label{def:l2}
Sea $\lambda>0$. La pérdida \emph{regularizada} con penalización $L^{2}$ es
\[
J_\lambda(\theta)=J(\theta)+\frac{\lambda}{2}\norm{\theta}^{2}.
\]
\end{definicion}

\begin{proposicion}[Encogimiento]\label{prop:encogimiento}
Un paso de descenso de gradiente sobre $J_\lambda$ es
\[
\theta_{k+1}=(1-\eta\lambda)\,\theta_k-\eta\,\nabla J(\theta_k).
\]
\end{proposicion}

Cada paso multiplica los pesos por $1-\eta\lambda<1$ antes de aplicar el gradiente: de ahí el nombre de \emph{decaimiento de pesos} (\emph{weight decay}).

\begin{proposicion}[Efecto sobre el minimizador]\label{prop:ridge}
Sea $J(\theta)=\frac12(\theta-c)\T Q(\theta-c)$ con $Q=\diag(a_1,\dots,a_n)$, $a_i>0$. El minimizador de $J_\lambda$ es
\[
\theta^{\ast}_i=\frac{a_i}{a_i+\lambda}\,c_i,\qquad i=1,\dots,n.
\]
\end{proposicion}

La regularización encoge hacia~0 cada coordenada en el factor $a_i/(a_i+\lambda)$: apenas afecta a las direcciones de curvatura grande ($a_i\gg\lambda$), en las que la pérdida es sensible, y anula casi por completo las de curvatura pequeña ($a_i\ll\lambda$)~\cite[§7.1.1]{gbc}. Por ejemplo, con $n=1$, $a=2$, $c=3$ y $\lambda=1$, el minimizador pasa de $3$ a $6/3=2$.

\begin{observacion}[Interpretación bayesiana]
Si $J$ es la log-verosimilitud negativa total de los datos y se adopta sobre $\theta$ una distribución a priori gaussiana $\mathcal{N}(0,\lambda^{-1}I)$, entonces $-\ln p(\theta)=\frac{\lambda}{2}\norm{\theta}^{2}+\text{cte}$, y minimizar $J_\lambda$ equivale a maximizar la distribución a posteriori (estimación MAP; módulo~5).
\end{observacion}

\section{Decaimiento de pesos desacoplado: AdamW}

Para el SGD, regularizar con $L^{2}$ y encoger los pesos en cada paso es lo mismo (\cref{prop:encogimiento}). Para un método preacondicionado no lo es~\cite[§2]{adamw}.

\begin{observacion}[Adam con $L^{2}$]\label{obs:adam-l2}
Si se aplica Adam a $J_\lambda$, el término $\lambda\theta$ se suma al gradiente antes de calcular $m_t$ y $v_t$, y queda dividido por $\sqrt{\hat v_t}+\varepsilon$. Cuando la normalización está dominada por $\nabla J$, la contribución del término de regularización al paso de la coordenada $i$ es aproximadamente
\[
\frac{\eta\lambda\,\theta_i}{\sqrt{\hat v_{t,i}}+\varepsilon},
\]
que depende de la escala del gradiente de esa coordenada, mientras que en el SGD es $\eta\lambda\theta_i$ para todas.
\end{observacion}

Las coordenadas con gradientes grandes ($\hat v_{t,i}$ grande) apenas se regularizan, y las de gradientes pequeños se regularizan en exceso. Loshchilov y Hutter~\cite{adamw} propusieron aplicar el decaimiento fuera de la normalización.

\begin{definicion}[AdamW]\label{def:adamw}
Con las cantidades $\hat m_t$ y $\hat v_t$ de la \cref{def:adam}, calculadas con el gradiente de $J$ sin regularizar, \emph{AdamW}~\cite[algoritmo~2]{adamw} realiza
\[
\theta_t=\theta_{t-1}-\eta_t\Bigl(\frac{\hat m_t}{\sqrt{\hat v_t}+\varepsilon}+\lambda\,\theta_{t-1}\Bigr),
\]
donde $\eta_t$ es la tasa de aprendizaje del calendario y $\lambda$ el coeficiente de decaimiento de pesos.
\end{definicion}

En AdamW todos los pesos se encogen en el mismo factor $1-\eta_t\lambda$, con independencia de sus gradientes. El \cref{tab:adamw} compara ambos métodos con $\eta=10^{-3}$, $\lambda=0{,}1$ y $\theta=1$.

\begin{table}[h]
\centering\small
\caption{Encogimiento de un peso $\theta=1$ en un paso, con $\eta=10^{-3}$ y $\lambda=0{,}1$.}\label{tab:adamw}
\begin{tabular}{@{}ccc@{}}
\toprule
$\hat v$ de la coordenada & Adam con $L^{2}$ ($\eta\lambda\theta/\sqrt{\hat v}$) & AdamW ($\eta\lambda\theta$) \\
\midrule
$10^{-4}$ & $0{,}01$ & $0{,}0001$ \\
$1$ & $0{,}0001$ & $0{,}0001$ \\
\bottomrule
\end{tabular}
\end{table}

\begin{observacion}[Convenios]
En~\cite{adamw}, el decaimiento de pesos del SGD se escribe $\theta_{t+1}=(1-\lambda')\theta_t-\alpha\nabla J(\theta_t)$, sin multiplicar $\lambda'$ por la tasa $\alpha$; con ese convenio, la equivalencia con $L^{2}$ exige $\lambda=\lambda'/\alpha$~\cite[proposición~1]{adamw}. Aquí se sigue el convenio de la \cref{def:adamw}, en el que el decaimiento por paso es $\eta_t\lambda$ y sigue al calendario. La proposición~2 de~\cite{adamw} demuestra que, para un preacondicionador que no sea múltiplo de la identidad, no existe ningún coeficiente $L^{2}$ que reproduzca el decaimiento desacoplado.
\end{observacion}

\addtocounter{definicion}{1}%

\section{Calendarios de la tasa de aprendizaje}

La tasa de aprendizaje es el hiperparámetro más influyente: si es demasiado grande, el método diverge (\cref{teo:gd-cuadratica}); si es demasiado pequeña, avanza despacio. Además, con gradientes estocásticos, un paso constante deja un suelo de ruido proporcional a $\eta$ (\cref{prop:suelo}). Por eso no se usa un $\eta$ fijo, sino un \emph{calendario} $\eta_t$~\cite[§12.11]{d2l}.

\begin{definicion}[Calentamiento lineal y decaimiento coseno]\label{def:coseno}
Sean $0\le\eta_{\min}<\eta_{\max}$, un número de pasos de calentamiento $W\ge 1$ y un número total de pasos $T>W$. El calendario de \emph{calentamiento lineal y decaimiento coseno} es
\[
\eta(t)=\begin{cases}
\eta_{\max}\,\dfrac{t}{W}, & 0\le t<W,\\[3mm]
\eta_{\min}+\dfrac{\eta_{\max}-\eta_{\min}}{2}\Bigl(1+\cos\Bigl(\pi\,\dfrac{t-W}{T-W}\Bigr)\Bigr), & W\le t\le T.
\end{cases}
\]
\end{definicion}

\begin{proposicion}[Propiedades]\label{prop:coseno}
El calendario de la \cref{def:coseno} es continuo en $t=W$, vale $\eta_{\max}$ en $t=W$ y $\eta_{\min}$ en $t=T$, es estrictamente decreciente en $[W,T]$ y en el punto medio del tramo coseno vale $(\eta_{\max}+\eta_{\min})/2$.
\end{proposicion}

\addtocounter{figure}{1}%

\addtocounter{definicion}{1}%

\paragraph{Calentamiento.} Al principio del entrenamiento, los gradientes son grandes e irregulares y los momentos de Adam se han estimado con muy pocos datos. Una tasa grande desde el primer paso puede hacer divergir la pérdida. Subir $\eta$ gradualmente da tiempo a que $\hat m_t$ y $\hat v_t$ se estabilicen; el calentamiento se popularizó para entrenar con lotes grandes~\cite{goyal}. El \cref{eje:5-sin-correccion} muestra que la propia corrección de sesgo de Adam actúa como un calentamiento implícito.

\paragraph{Decaimiento.} Por la \cref{prop:suelo}, el exceso de pérdida estacionario del SGD es proporcional a $\eta$; reducir $\eta$ al final permite que los parámetros se asienten. El decaimiento coseno procede del método SGDR de Loshchilov y Hutter~\cite{sgdr}, que lo combinaba con reinicios periódicos; en el preentrenamiento se usa un único ciclo.

\begin{observacion}[La receta de los modelos de lenguaje abiertos]
LLaMA documenta exactamente esta receta: AdamW con $\beta_1=0{,}9$, $\beta_2=0{,}95$, decaimiento de pesos $\lambda=0{,}1$, recorte del gradiente a norma~1, 2000 pasos de calentamiento y decaimiento coseno hasta el 10~\% de la tasa máxima~\cite[§2.3]{llama}. Es la misma estructura que la del \cref{ej:coseno}. Como se indicó en el \cref{tema:5}, Anthropic, Google y xAI no han publicado los calendarios ni los optimizadores de Claude, Gemini y Grok.
\end{observacion}



\backmatter
% =====================================================================
\begin{thebibliography}{99}
\addcontentsline{toc}{chapter}{Bibliografía}
\markboth{Bibliografía}{}
\small
\def\UrlBreaks{\do\/\do-}

\bibitem{bottou}
L.~Bottou, F.~E. Curtis y J.~Nocedal.
\newblock Optimization methods for large-scale machine learning.
\newblock \emph{SIAM Review}, 60(2):223--311, 2018.
\newblock \url{https://arxiv.org/abs/1606.04838}

\bibitem{bv}
S.~Boyd y L.~Vandenberghe.
\newblock \emph{Convex Optimization}.
\newblock Cambridge University Press, 2004.
\newblock Cap.~2, «Convex sets»; cap.~3, «Convex functions»; cap.~9, «Unconstrained minimization». Texto del curso EE364a de Stanford.
\newblock \url{https://web.stanford.edu/~boyd/cvxbook/}

\bibitem{ee364a}
S.~Boyd.
\newblock EE364a: Convex Optimization I.
\newblock Stanford University. Transparencias del curso.
\newblock \url{https://web.stanford.edu/class/ee364a/lectures.html}

\bibitem{gpt3}
T.~B. Brown \emph{et al.}
\newblock Language models are few-shot learners.
\newblock En \emph{Advances in Neural Information Processing Systems 33 (NeurIPS)}, 2020.
\newblock \url{https://arxiv.org/abs/2005.14165}

\bibitem{cohen}
J.~M. Cohen, S.~Kaur, Y.~Li, J.~Z. Kolter y A.~Talwalkar.
\newblock Gradient descent on neural networks typically occurs at the edge of stability.
\newblock En \emph{International Conference on Learning Representations (ICLR)}, 2021.
\newblock \url{https://arxiv.org/abs/2103.00065}

\bibitem{dauphin}
Y.~N. Dauphin, R.~Pascanu, C.~Gulcehre, K.~Cho, S.~Ganguli e Y.~Bengio.
\newblock Identifying and attacking the saddle point problem in high-dimensional non-convex optimization.
\newblock En \emph{Advances in Neural Information Processing Systems 27 (NeurIPS)}, 2014.
\newblock \url{https://arxiv.org/abs/1406.2572}

\bibitem{adagrad}
J.~Duchi, E.~Hazan e Y.~Singer.
\newblock Adaptive subgradient methods for online learning and stochastic optimization.
\newblock \emph{Journal of Machine Learning Research}, 12:2121--2159, 2011.

\bibitem{goh}
G.~Goh.
\newblock Why momentum really works.
\newblock \emph{Distill}, 2017.
\newblock \url{https://distill.pub/2017/momentum/}

\bibitem{gbc}
I.~Goodfellow, Y.~Bengio y A.~Courville.
\newblock \emph{Deep Learning}.
\newblock MIT Press, 2016.
\newblock Cap.~4, «Numerical Computation»; cap.~7, «Regularization for Deep Learning»; cap.~8, «Optimization for Training Deep Models».
\newblock \url{https://www.deeplearningbook.org/}

\bibitem{goyal}
P.~Goyal, P.~Dollár, R.~Girshick \emph{et al.}
\newblock Accurate, large minibatch SGD: Training ImageNet in 1 hour.
\newblock Prepublicación, 2017.
\newblock \url{https://arxiv.org/abs/1706.02677}

\bibitem{rmsprop}
G.~Hinton, N.~Srivastava y K.~Swersky.
\newblock Neural Networks for Machine Learning, clase~6e: «rmsprop: Divide the gradient by a running average of its recent magnitude».
\newblock Curso en línea, Coursera, 2012.

\bibitem{keskar}
N.~S. Keskar, D.~Mudigere, J.~Nocedal, M.~Smelyanskiy y P.~T.~P. Tang.
\newblock On large-batch training for deep learning: Generalization gap and sharp minima.
\newblock En \emph{International Conference on Learning Representations (ICLR)}, 2017.
\newblock \url{https://arxiv.org/abs/1609.04836}

\bibitem{adam}
D.~P. Kingma y J.~Ba.
\newblock Adam: A method for stochastic optimization.
\newblock En \emph{International Conference on Learning Representations (ICLR)}, 2015.
\newblock \url{https://arxiv.org/abs/1412.6980}

\bibitem{li-paisaje}
H.~Li, Z.~Xu, G.~Taylor, C.~Studer y T.~Goldstein.
\newblock Visualizing the loss landscape of neural nets.
\newblock En \emph{Advances in Neural Information Processing Systems 31 (NeurIPS)}, 2018.
\newblock \url{https://arxiv.org/abs/1712.09913}

\bibitem{sgdr}
I.~Loshchilov y F.~Hutter.
\newblock SGDR: Stochastic gradient descent with warm restarts.
\newblock En \emph{International Conference on Learning Representations (ICLR)}, 2017.
\newblock \url{https://arxiv.org/abs/1608.03983}

\bibitem{adamw}
I.~Loshchilov y F.~Hutter.
\newblock Decoupled weight decay regularization.
\newblock En \emph{International Conference on Learning Representations (ICLR)}, 2019.
\newblock \url{https://arxiv.org/abs/1711.05101}

\bibitem{nesterov}
Y.~Nesterov.
\newblock A method for solving the convex programming problem with convergence rate $O(1/k^{2})$.
\newblock \emph{Soviet Mathematics Doklady}, 27:372--376, 1983.

\bibitem{polyak}
B.~T. Polyak.
\newblock Some methods of speeding up the convergence of iteration methods.
\newblock \emph{USSR Computational Mathematics and Mathematical Physics}, 4(5):1--17, 1964.

\bibitem{robbins}
H.~Robbins y S.~Monro.
\newblock A stochastic approximation method.
\newblock \emph{The Annals of Mathematical Statistics}, 22(3):400--407, 1951.

\bibitem{sagun}
L.~Sagun, L.~Bottou e Y.~LeCun.
\newblock Eigenvalues of the Hessian in deep learning: Singularity and beyond.
\newblock Prepublicación, 2016.
\newblock \url{https://arxiv.org/abs/1611.07476}

\bibitem{strang}
G.~Strang.
\newblock \emph{Linear Algebra and Learning from Data}.
\newblock Wellesley-Cambridge Press, 2019.
\newblock Texto del curso 18.065 del MIT.

\bibitem{mit18065}
G.~Strang.
\newblock 18.065: Matrix Methods in Data Analysis, Signal Processing, and Machine Learning.
\newblock Massachusetts Institute of Technology, MIT OpenCourseWare, primavera de 2018.
\newblock \url{https://ocw.mit.edu/courses/18-065-matrix-methods-in-data-analysis-signal-processing-and-machine-learning-spring-2018/}

\bibitem{llama}
H.~Touvron \emph{et al.}
\newblock LLaMA: Open and efficient foundation language models.
\newblock Prepublicación, 2023.
\newblock \url{https://arxiv.org/abs/2302.13971}

\bibitem{d2l}
A.~Zhang, Z.~C. Lipton, M.~Li y A.~J. Smola.
\newblock \emph{Dive into Deep Learning}, versión 1.0.3.
\newblock Cambridge University Press, 2023.
\newblock Cap.~12, «Optimization Algorithms».
\newblock \url{https://d2l.ai/}

\end{thebibliography}


\end{document}
